\subsection{Inclusions of function spaces in the space of Schwartz functions}

PROPOSITION 12. — The space \(\mathcal{D}(\mathbf{R}^{n})\) is contained in \(\mathcal{S}(\mathbf{R}^{n})\) , and the inclusion of \(\mathcal{D}(\mathbf{R}^{n})\) into \(\mathcal{S}(\mathbf{R}^{n})\) is continuous with dense image.

Let \(B \subset \mathcal{D}(\mathbf{R}^{n})\) be a bounded subset, and let K be a compact subset of \(\mathbf{R}^{n}\) such that \(\mathrm{B} \subset \mathcal{C}_{\mathrm{K}}^{\infty}(\mathbf{R}^{n})\) . Let \(\alpha \in \mathbf{N}^{n}\) and \(k \in \mathbf{N}\) . For every function \(\varphi \in B\) , it follows that

\[
\widetilde {q} _ {\alpha , k} (\varphi) \leqslant \Bigl (\sup _ {x \in \mathrm{K}} \| x \| ^ {k} \Bigr) p _ {\alpha , \mathrm{K}} (\varphi),
\]

so B is bounded in \(\mathcal{S}(\mathbf{R}^{n})\) . The continuity of the inclusion then follows from the fact that the spaces \(\mathcal{S}(\mathbf{R}^{n})\) and \(\mathcal{D}(\mathbf{R}^{n})\) are bornological.

Let us prove that \(\mathcal{D}(\mathbf{R}^n)\) is dense in \(\mathcal{S}(\mathbf{R}^n)\) . Let B be the unit ball of \(\mathbf{R}^n\) and \(\eta \in \mathcal{D}(\mathbf{R}^n)\) a test function with support contained in 2B such that \(0 \leqslant \eta \leqslant 1\) and \(\eta(x) = 1\) for all \(x \in \mathrm{B}\) (lemma 1, \(a\) ) of IV, p. 196).

Let \(\varphi \in \mathcal{S}(\mathbf{R}^n)\) . For every integer \(m \geqslant 1\) and every \(x \in \mathbf{R}^n\) , set \(\eta_m(x) = \eta(x/m)\) . Finally define \(\varphi_m = \eta_m\varphi\) ; we have \(\varphi_m \in \mathcal{D}(\mathbf{R}^n)\) .

Since \(\partial^{\alpha}\eta_{m} = m^{-|\alpha|}(\partial^{\alpha}\eta)(x / m)\) for all \(\alpha \in \mathbf{N}^{n}\) and \(x\in \mathbf{R}^n\) , from formula (2) of IV, p. 210 we deduce that the sequence \((\varphi_{m})\) is bounded in \(\mathcal{S}(\mathbf{R}^n)\) .

Let C be a compact subset of \(R^{n}\) . The sequence \((\varphi_{m})_{m\geqslant1}\) converges to \(\varphi\) in \(\mathcal{C}_{\mathrm{K}}^{\infty}(\mathbf{R}^{n})\) since \(\varphi_{m}\) coincides with \(\varphi\) on C for all sufficiently large m. Thus the sequence \((\varphi_{m})\) converges to \(\varphi\) in \(\mathcal{C}^{\infty}(\mathbf{R}^{n})\) , and the corollary of Proposition 10 of IV, p. 211 allows us to conclude that the sequence \((\varphi_{m})\) converges to \(\varphi\) in \(\mathcal{S}(\mathbf{R}^{n})\) .

Lemma 8. --- Let B be a bounded subset of \(\mathcal{D}(\mathbf{R}^{n})\) . The topology induced on B by the topology of \(\mathcal{S}(\mathbf{R}^{n})\) coincides with the topology induced by \(\mathcal{D}(\mathbf{R}^{n})\) .

Since the inclusion of \(\mathcal{D}(\mathbf{R}^{n})\) into \(\mathcal{S}(\mathbf{R}^{n})\) is continuous, the topology on B induced by \(\mathcal{D}(\mathbf{R}^{n})\) is finer than the one induced by \(\mathcal{S}(\mathbf{R}^{n})\) . Moreover, there exists a compact subset K of \(\mathbf{R}^{n}\) such that \(\mathrm{B} \subset \mathcal{C}_{\mathrm{K}}^{\infty}(\mathbf{R}^{n})\) . For every \(\alpha \in \mathbf{N}^{n}\) , we then have \(p_{\alpha,K}(\varphi) \leqslant \widetilde{q}_{\alpha,0}(\varphi)\) , which implies that the topology induced by \(\mathcal{S}(\mathbf{R}^{n})\) is finer than the one induced by \(\mathcal{D}(\mathbf{R}^{n})\) .

PROPOSITION 13. --- Let \(p \in [1, +\infty]\) . The space \(\mathcal{S}(\mathbf{R}^{n})\) is contained in \(\mathcal{L}^{p}(\mathbf{R}^{n})\) and the canonical injection of \(\mathcal{S}(\mathbf{R}^{n})\) into \(\mathcal{L}^{p}(\mathbf{R}^{n})\) is continuous. The image of \(\mathcal{S}(\mathbf{R}^{n})\) in \(\mathrm{L}^{p}(\mathbf{R}^{n})\) is dense if \(p \neq +\infty\) .

The first assertion is immediate for \(p = +\infty\) . Suppose now that \(p \in [1, +\infty[\) . Let m be an integer such that \(n + 1 < mp\) . For every \(\varphi \in \mathcal{S}(\mathbf{R}^{n})\) and \(x \in \mathbf{R}^{n}\) , we have

\[
\| x \| ^ {n + 1} | \varphi (x) | ^ {p} \leqslant \widetilde {q} _ {0, m} (\varphi) ^ {p}
\]

hence \(\varphi \in \mathcal{L}^p (\mathbf{R}^n)\) by prop. 3 of IV, p. 199. Moreover, we obtain

\[
\mathrm{N} _ {p} (\varphi) \leqslant a _ {n} ^ {1 / p} \widetilde {q} _ {0, 0} (\varphi) + b _ {n} \widetilde {q} _ {0, m} (\varphi),
\]

where

\[
a _ {n} = \int_ {\| x \| \leqslant 1} d \mu (x), \quad b _ {n} = \int_ {\| x \| \geqslant 1} \frac {1}{\| x \| ^ {n + 1}} d \mu (x),
\]

therefore the injection of \(\mathcal{S}(\mathbf{R}^n)\) into \(\mathcal{L}^p (\mathbf{R}^n)\) is continuous.

As the space \(\mathcal{D}(\mathbf{R}^n)\) is contained in \(\mathcal{S}(\mathbf{R}^n)\) , Proposition 4 of IV, p. 202 implies that \(\mathcal{S}(\mathbf{R}^n)\) is dense in \(\mathrm{L}^p (\mathbf{R}^n)\) if \(p\neq +\infty\) .

